\documentclass[11pt]{article}

\usepackage[margin=1in]{geometry}
\usepackage{lmodern}
\usepackage[T1]{fontenc}
\usepackage{microtype}
\usepackage{amsmath}
\usepackage{parskip}

\begin{document}

\begin{flushright}
  \today
\end{flushright}

\vspace{1em}

\noindent
Editors of \textit{Mathematics}

\vspace{1em}

\noindent
Dear Editors,

We are pleased to submit our manuscript, entitled
\textit{``The Missing Signal: Gradient Injection Recovers $O(n^2)$ Residual
Decoding in GNNs,''} for consideration as an Article in \textit{Mathematics}.

The manuscript studies a concrete limitation of neural decoders for combinatorial
optimization. Classical residual algorithms are efficient because they maintain
statistics, such as degree, that can be updated locally after each deletion.
GNN-based scores, however, pass through nonlinear transformations and generally do
not admit such local additive updates. Consequently, a learned residual decoder may
need to rerun the full GNN after every deletion, increasing the decoding cost from
$O(n^2)$ to $O(n^3)$.

Our contribution is to identify this update-compatibility issue and show that, for
planted-clique recovery, the probability-space gradient of an unsupervised clique
objective provides a valid residual coordinate. The key point is not simply that
gradients are useful auxiliary features, but that this gradient has an exact local
update rule under vertex deletion. Its coordinates can therefore be maintained in
$O(n)$ time per deletion, allowing the decoder to run the GNN once and then update a
cached residual state throughout the pruning process.

Empirically, the resulting SGNN+GU decoder with Self-Imitation Training matches the
rerun-based $O(n^3)$ decoder on Easy and Medium planted-clique regimes while
retaining $O(n^2)$ decoding complexity.

We believe the manuscript is suitable for \textit{Mathematics} because it connects
graph algorithms, combinatorial optimization, planted-clique recovery, and the
mathematical analysis of neural algorithmic reasoning. The central question is not
only whether a GNN can learn a useful score, but whether the quantities used by a
learned decoder have valid residual dynamics that support efficient algorithmic
updates.

We confirm that neither the manuscript nor any part of its content is currently
under consideration for publication in another journal, and that the manuscript has
not been previously published. All authors have approved the manuscript and agree
with its submission to \textit{Mathematics}.

\vspace{2em}

\noindent
Sincerely,

\vspace{1.5em}

\noindent
Elad Shoham\\
Department of Software and Information Systems Engineering\\
Ben-Gurion University of the Negev\\
Beer-Sheva 8410501, Israel\\
\texttt{shellad@post.bgu.ac.il}

\vspace{1em}

\noindent
On behalf of all authors:\\
Elad Shoham, Havana Rika, and Dan Vilenchik

\end{document}
